\documentclass[11pt]{article}

\usepackage[margin=1in]{geometry}
\usepackage{booktabs}
\usepackage{amsmath}
\usepackage{enumitem}
\usepackage{hyperref}

\title{Design of Mahimahi-Inspired Congestion-Control Environments}
\author{}
\date{}

\begin{document}
\maketitle

\section{Overview}

The repository defines four scenario-shaped Gymnasium environments for
congestion-control experiments:

\begin{enumerate}[label=(\arabic*)]
  \item a stable wired bottleneck,
  \item a cellular trace with bursty capacity variation,
  \item a LEO-inspired trace with abrupt delay and bandwidth changes, and
  \item a multi-flow fairness stress test.
\end{enumerate}

Each scenario defaults to a fast in-process simulator, but can also be run
through the real \texttt{mm-link} subprocess adapter by passing
\texttt{use\_mahimahi=True}. In simulator mode, the environments model the same
core signals that Mahimahi-style experiments expose: bandwidth, propagation
delay, queueing delay, packet loss, throughput, and multi-flow sharing. In real
Mahimahi mode, the scenario classes delegate to scenario-shaped wrappers around
\texttt{MahimahiSubprocessEnv}, preserving the same action and observation
interfaces while launching \texttt{mm-link} with scenario-specific trace files.

\section{Shared Single-Flow Design}

The wired, cellular, and LEO environments share a single-flow bottleneck model
in simulator mode. At each step, the agent chooses a target sending rate:

\[
  a_t = r_t \quad \text{in Mbps.}
\]

The environment replays a link trace with a capacity \(C_t\), propagation delay
\(D_t\), and fixed step duration. If the selected sending rate exceeds the
available service rate, data accumulates in a finite bottleneck queue. Once the
queue is full, excess traffic is dropped.

The single-flow observation is seven-dimensional:

\[
o_t =
\left[
\frac{r_t}{r_{\max}},
\frac{x_t}{r_{\max}},
\frac{C_t}{r_{\max}},
\frac{D_t}{D_{\max}},
\frac{Q_t}{Q_{\max}},
L_t,
B_t
\right],
\]

where \(r_t\) is the chosen send rate, \(x_t\) is delivered throughput, \(C_t\)
is link capacity, \(D_t\) is propagation delay, \(Q_t\) is queueing delay,
\(L_t\) is loss fraction, and \(B_t\) is normalized queue occupancy.

The reward is:

\[
R_t = x_t - \lambda_D (D_t + Q_t) - \lambda_L L_t.
\]

This reward encourages high delivered throughput while penalizing round-trip
latency and loss. The exact delay and loss penalties can vary by scenario. The
default single-flow loss penalty is \(\lambda_L = 10.0\). The default delay
penalty is \(\lambda_D = 0.015\), except for the LEO-inspired scenario, which
lowers it to \(\lambda_D = 0.012\).

\subsection{Optional MPCC-Style Reward}

The scenarios also support an optional MPCC-style reward by setting
\texttt{reward\_mode=mpcc}. The default reward above remains available as
\texttt{reward\_mode=default}; MPCC mode changes only the scalar reward and
diagnostic fields in \texttt{info}, not the action or observation spaces.

The MPCC reward uses smoothed measurements \(\hat{T}\) and \(\hat{R}\) for
throughput and RTT. The simulator updates these smoothed values with first-order
filters using configurable time constants \(\tau_T\) and \(\tau_R\). In
single-flow simulator mode, \(R_0\) is the current propagation delay. In the
multi-flow simulator, which has no explicit propagation-delay trace, \(R_0\)
defaults to \(20\) ms and can be overridden with
\texttt{mpcc\_base\_rtt\_ms}.

The reference curve in throughput--RTT space is:

\[
\Gamma(\theta) =
\left(C\theta,\; R_0 + \alpha\theta^2\right),
\quad
\theta = \mathrm{clip}\left(\frac{\hat{T}}{C}, 0, 1\right).
\]

The unit tangent is:

\[
\mathbf{t}(\theta) =
\frac{(C,\;2\alpha\theta)}
{\sqrt{C^2 + 4\alpha^2\theta^2}}
= (t_T,t_R).
\]

Given \(\Gamma(\theta)=(\Gamma_T,\Gamma_R)\), the contouring and lag errors are:

\[
e_c = t_R(\hat{T}-\Gamma_T) - t_T(\hat{R}-\Gamma_R),
\]

\[
e_\ell = t_T(\hat{T}-\Gamma_T) + t_R(\hat{R}-\Gamma_R).
\]

The MPCC stage cost implemented by the environments is:

\[
\ell_{\mathrm{mpcc}} =
w_c e_c^2
+ w_\ell e_\ell^2
+ w_d [\hat{R}-R^*]_+^2
+ w_u \frac{\sum_i (s_{i,t}-s_{i,t-1})^2}{s_{\max}^2}
+ w_{\mathrm{loss}} L_t^2
- w_p
\frac{\hat{T}/C}{\hat{R}/R_0}.
\]

Gymnasium maximizes reward, so the environment returns:

\[
R_t^{\mathrm{mpcc}} = -\ell_{\mathrm{mpcc}}.
\]

The default MPCC weights are intentionally modest because the geometric errors
are measured in physical units:

\[
w_c=10^{-3},\quad
w_\ell=10^{-3},\quad
w_d=10^{-4},\quad
w_p=1.0,\quad
w_u=0.1,\quad
w_{\mathrm{loss}}=0.
\]

If \(\alpha\) is not provided, it defaults to half the queue-delay scale. If
\(R^*\) is not provided, it defaults to \(R_0\) plus one quarter of the
queue-delay scale.

The \texttt{info} dictionary exposes MPCC terms for debugging and analysis:

\begin{itemize}
  \item \url{mpcc_stage_cost}, \url{mpcc_contour_error}, and
        \url{mpcc_lag_error};
  \item \url{mpcc_power}, \url{mpcc_smoothing_penalty}, and
        \url{mpcc_theta};
  \item \url{mpcc_smoothed_throughput_mbps} and
        \url{mpcc_smoothed_rtt_ms}.
\end{itemize}

\subsection{Real Mahimahi Single-Flow Mode}

When \texttt{use\_mahimahi=True}, the single-flow scenarios instantiate
\texttt{SingleFlowMahimahiScenarioEnv}. This wrapper delegates stepping to
\texttt{MahimahiSubprocessEnv} with one action dimension, then maps the
adapter's measurement dictionary back into the seven-dimensional scenario
observation. Since a real controller may report queueing delay without explicit
queue occupancy, the final queue-related observation slot is represented by
normalized queueing delay in this mode.

The default controller command is:

\begin{verbatim}
python examples/mm_link_json_controller.py
\end{verbatim}

The wrapper forwards \texttt{controller\_cmd}, \texttt{mahimahi\_cmd},
\texttt{timeout\_s}, \texttt{restart\_on\_reset}, and \texttt{render\_mode} to
the subprocess adapter. In \texttt{reward\_mode=mpcc}, the wrapper replaces the
adapter's default reward with the same MPCC reward calculation using the
measurements reported by the controller.

\section{Stable Wired Bottleneck}

\textbf{Environment ID:} \texttt{MahiWiredBottleneck-v0}

\textbf{Class:} \texttt{WiredBottleneckEnv}

\subsection{Design Goal}

This environment is a stationary sanity-check task. The optimal behavior should
be easy to understand: send near the fixed bottleneck rate, avoid persistent
oversending, and keep delay/loss low.

\subsection{Configuration}

\begin{itemize}
  \item Capacity: \(10\) Mbps
  \item Propagation delay: \(20\) ms
  \item Step duration: \(100\) ms
  \item Queue size: \(150\) ms
  \item Maximum send rate: \(20\) Mbps
  \item Real Mahimahi traces: \texttt{mahimahi/12mbps.trace} for both uplink
        and downlink
\end{itemize}

\subsection{Rationale}

A fixed \(10\) Mbps link creates a clean, stationary bottleneck. The low
\(20\) ms propagation delay resembles a wired path rather than a cellular or
satellite path. The queue is intentionally modest: large enough to permit
transient overshoot, but small enough that persistent oversending produces
visible queueing delay and loss. The \(20\) Mbps maximum action lets the agent
overshoot by a factor of two, so it can learn the cost of aggression.

\subsection{Expected Policy}

A good policy should converge to sending close to \(10\) Mbps, with low loss and
low queueing delay.

\section{Cellular Bursty Capacity}

\textbf{Environment ID:} \texttt{MahiCellularBursty-v0}

\textbf{Class:} \texttt{CellularBurstyEnv}

\subsection{Design Goal}

This environment represents a cellular link where radio conditions cause sharp
capacity variation. The agent must adapt to bursts and collapses without filling
the queue.

\subsection{Capacity Trace}

\[
[8.0, 7.5, 3.0, 1.2, 2.0, 12.0, 15.0, 6.0, 4.0, 2.2,
1.0, 9.0, 14.0, 18.0, 5.0, 2.5, 1.5, 7.0, 11.0, 4.5]
\]

\subsection{Delay Trace}

\[
[55, 58, 70, 95, 85, 60, 50, 65, 78, 90,
110, 62, 54, 48, 72, 88, 105, 68, 58, 76]
\]

\subsection{Rationale}

Capacity ranges from roughly \(1\) Mbps to \(18\) Mbps. Low-capacity periods
punish aggressive sending, while high-capacity bursts reward rapid utilization.
The delay trace rises during weaker-capacity periods to mimic cellular
instability and buffer pressure. The queue size is \(300\) ms, larger than the
wired queue, because cellular paths often exhibit deeper and more variable
buffering.

In real Mahimahi mode, this scenario uses
\texttt{mahimahi/traces/TMobile-LTE-driving.up} as the default uplink trace and
\texttt{mahimahi/traces/TMobile-LTE-driving.down} as the default downlink trace.
The default maximum send rate is \(30\) Mbps.

\subsection{Expected Policy}

A good policy should increase its sending rate during high-capacity bursts and
back off during collapses, while avoiding persistent cellular queue buildup.

\section{LEO-Inspired Satellite Trace}

\textbf{Environment ID:} \texttt{MahiLeoSatellite-v0}

\textbf{Class:} \texttt{LeoSatelliteEnv}

\subsection{Design Goal}

This environment models abrupt bandwidth and propagation-delay shifts inspired
by low-Earth-orbit satellite networking. It is not a detailed orbital model;
the goal is to stress robustness to sudden regime changes.

\subsection{Capacity Trace}

\[
[35, 35, 32, 8, 6, 45, 48, 12, 10, 9,
50, 52, 18, 15, 4, 4, 30, 42, 46, 7]
\]

\subsection{Delay Trace}

\[
[35, 36, 38, 80, 95, 42, 36, 75, 88, 92,
40, 34, 65, 70, 120, 125, 55, 44, 38, 100]
\]

\subsection{Rationale}

The capacity trace includes high-throughput windows near \(50\) Mbps and sudden
collapses to \(4\)--\(8\) Mbps. The delay trace includes abrupt jumps from
roughly \(35\) ms to more than \(120\) ms. These shifts are meant to mimic the
kind of behavior that can arise from satellite handoff, path geometry changes,
obstruction, or beam transitions.

The delay penalty is slightly lower than in some other scenarios:

\[
\lambda_D = 0.012.
\]

This prevents the policy from becoming overly conservative in a setting where
propagation delay variation is inherent.

In real Mahimahi mode, this scenario uses
\texttt{mahimahi/traces/Verizon-LTE-short.up} as the default uplink trace and
\texttt{mahimahi/traces/Verizon-LTE-short.down} as the default downlink trace.
The default maximum send rate is \(60\) Mbps.

\subsection{Expected Policy}

A good policy should exploit high-capacity windows while remaining robust to
sudden bandwidth drops and delay jumps.

\section{Multi-Flow Fairness Stress Test}

\textbf{Environment ID:} \texttt{MahiMultiFlowFairness-v0}

\textbf{Class:} \texttt{MultiFlowFairnessEnv}

\subsection{Design Goal}

This environment tests whether a controller can balance total throughput and
fairness when multiple flows share a single bottleneck.

\subsection{Default Configuration}

\begin{itemize}
  \item Number of flows: \(3\)
  \item Shared bottleneck capacity: \(12\) Mbps
  \item Maximum send rate per flow: \(15\) Mbps
  \item Queue size: \(200\) ms
  \item Real Mahimahi traces: \texttt{mahimahi/12mbps.trace} for both uplink
        and downlink
\end{itemize}

\subsection{Action}

The action is a vector of per-flow sending rates:

\[
a_t = [r_{1,t}, r_{2,t}, r_{3,t}].
\]

\subsection{Observation}

For \(N\) flows, the observation is:

\[
o_t =
\left[
\frac{r_{1,t}}{r_{\max}}, \ldots, \frac{r_{N,t}}{r_{\max}},
\frac{x_{1,t}}{r_{\max}}, \ldots, \frac{x_{N,t}}{r_{\max}},
\frac{C}{r_{\max}},
B_t,
L_t,
J_t
\right],
\]

where \(J_t\) is Jain's fairness index.

In simulator mode, \(B_t\) is normalized queue occupancy. In real Mahimahi
mode, the corresponding slot is normalized queueing delay, because the
subprocess adapter may not receive an explicit queue occupancy measurement from
the controller.

\subsection{Reward}

The reward combines total throughput, fairness, delay, and loss:

\[
R_t =
\sum_{i=1}^{N} x_{i,t}
+ \lambda_F J_t
- \lambda_D Q_t
- \lambda_L L_t.
\]

The default fairness weight is:

\[
\lambda_F = 8.0.
\]

The default delay and loss penalties are \(\lambda_D = 0.015\) and
\(\lambda_L = 4.0\).

Jain's fairness index is:

\[
J_t =
\frac{\left(\sum_i x_{i,t}\right)^2}
{N \sum_i x_{i,t}^2}.
\]

In \texttt{reward\_mode=mpcc}, the multi-flow environment applies the same
throughput--RTT MPCC objective to aggregate throughput and shared queueing
delay. It additionally adds a fair-share penalty:

\[
\ell_{\mathrm{fair}} =
w_f \sum_{i=1}^{N}
\left(s_{i,t} - \frac{C}{N}\right)^2.
\]

This term biases each sender toward its nominal share of the bottleneck
capacity. The default MPCC fair-share weight is:

\[
w_f = 10^{-2}.
\]

The total MPCC reward for the multi-flow scenario is:

\[
R_t^{\mathrm{mpcc,multi}} =
-\left(\ell_{\mathrm{mpcc}} + \ell_{\mathrm{fair}}\right),
\]

where \(\ell_{\mathrm{mpcc}}\) includes the contouring, lag, delay, power,
smoothing, and optional loss terms defined in the shared MPCC reward section.
The \texttt{info} dictionary includes \texttt{mpcc\_fairness\_penalty} in this
mode.

\subsection{Rationale}

If the reward only used total throughput, an agent could learn unfair
allocations where one flow dominates. Jain's fairness term rewards equal
sharing. Since each flow can request up to \(15\) Mbps but the shared bottleneck
is only \(12\) Mbps, the combined action can dramatically oversubscribe the
link. This forces the agent to learn both rate control and fair allocation.

\subsection{Expected Policy}

A good policy should keep aggregate sending near the shared capacity, distribute
throughput roughly equally across flows, maintain high Jain fairness, and avoid
queue buildup and loss.

\subsection{Real Mahimahi Multi-Flow Mode}

When \texttt{use\_mahimahi=True}, \texttt{MultiFlowFairnessEnv} delegates to
\texttt{MultiFlowMahimahiScenarioEnv}. The wrapper configures
\texttt{MahimahiSubprocessEnv} with \texttt{action\_dim = N}, then reconstructs
per-flow throughput observations from the controller's
\texttt{throughputs\_mbps} field when available. If only aggregate throughput is
reported, the wrapper distributes it across flows in proportion to the chosen
send rates before computing Jain's fairness index. In
\texttt{reward\_mode=mpcc}, the wrapper computes the MPCC reward from the
reported aggregate or reconstructed per-flow throughput, queueing delay, RTT,
and capacity measurements.

\section{Comparison}

\begin{table}[h]
\centering
\begin{tabular}{lll}
\toprule
Environment & Main Challenge & Action Space \\
\midrule
\texttt{MahiWiredBottleneck-v0} & Stable bottleneck convergence & One send rate \\
\texttt{MahiCellularBursty-v0} & Bursty capacity adaptation & One send rate \\
\texttt{MahiLeoSatellite-v0} & Abrupt bandwidth/delay shifts & One send rate \\
\texttt{MahiMultiFlowFairness-v0} & Shared bottleneck fairness & Per-flow send rates \\
\bottomrule
\end{tabular}
\caption{Summary of the four scenario environments.}
\end{table}

\section{Real Mahimahi Adapter}

The scenario environments use \texttt{MahimahiSubprocessEnv} internally when
\texttt{use\_mahimahi=True}. The adapter connects to Mahimahi by launching:

\begin{verbatim}
mm-link uplink.trace downlink.trace -- <controller command>
\end{verbatim}

The Gym environment communicates with the controller process over JSON lines.
At each step, the environment writes an action to the controller's standard
input:

\begin{verbatim}
{"type": "action", "step": 0, "send_rates_mbps": [4.0]}
\end{verbatim}

The controller must print one measurement object to standard output:

\begin{verbatim}
{"throughput_mbps": 3.8, "rtt_ms": 45.0, "queue_delay_ms": 5.0,
 "loss_fraction": 0.0, "capacity_mbps": 10.0}
\end{verbatim}

This adapter is intentionally generic. The controller command can be a Python
shim, a compiled congestion-control binary, or a script that manages a real
sender and receiver inside the Mahimahi shell. It remains an internal building
block for the four scenario environments rather than a separately registered
Gymnasium environment.

The built-in scenario wrappers use the following default trace files for real
Mahimahi runs:

\begin{table}[h]
\centering
\small
\begin{tabular}{p{0.31\linewidth}p{0.29\linewidth}p{0.29\linewidth}}
\toprule
Environment & Uplink trace & Downlink trace \\
\midrule
\texttt{MahiWiredBottleneck-v0} & \url{mahimahi/12mbps.trace} & \url{mahimahi/12mbps.trace} \\
\texttt{MahiCellularBursty-v0} & \url{mahimahi/traces/TMobile-LTE-driving.up} & \url{mahimahi/traces/TMobile-LTE-driving.down} \\
\texttt{MahiLeoSatellite-v0} & \url{mahimahi/traces/Verizon-LTE-short.up} & \url{mahimahi/traces/Verizon-LTE-short.down} \\
\texttt{MahiMultiFlowFairness-v0} & \url{mahimahi/12mbps.trace} & \url{mahimahi/12mbps.trace} \\
\bottomrule
\end{tabular}
\caption{Default real Mahimahi trace files used by the scenario wrappers.}
\end{table}

Mahimahi trace reminder: \texttt{mm-link} traces are packet-delivery schedules,
where each line is a delivery opportunity timestamp for an MTU-sized packet.
They are not simple \texttt{duration\_ms capacity\_mbps} rows. The in-process
simulator traces use Python arrays or simulator-specific trace formats, while
real Mahimahi runs should use Mahimahi-compatible uplink and downlink trace
files.

\section{Scenario Configuration Direction}

A natural next step is to move the remaining hard-coded scenario parameters into
configuration files that can be consumed by both the simulator and the real
Mahimahi subprocess wrapper. For example:

\begin{verbatim}
name: cellular_bursty
capacity_trace: ...
delay_trace: ...
queue_size_ms: 300
flows: 1
reward:
  mode: mpcc
  contour_weight: 0.001
  lag_weight: 0.001
  delay_weight: 0.0001
  power_weight: 1.0
  smoothing_weight: 0.1
  fairness_weight: 0.01
\end{verbatim}

This would allow policies to be trained quickly in simulation and then evaluated
against real \texttt{mm-link} experiments with the same scenario definitions.

\end{document}
